\scanpage{71}
\uline{Morphismes essentiellement affines}

$f:X\longrightarrow Y$ est dit ess. affine s'il se factorise en $X\xrightarrow{\deleted{\unsure{g}}f_1}X_0\xrightarrow{f_0}Y$, avec $f_0$ de prés. finie et $f_1$ affine. Il revient au $\overline m$, \uline{si} $Y$ est \inserted{loc.} noethérien\deleted{\unsure{noequ.|ou qu.}}, que $X$ est isomorphe à la $\varprojlim$ d'un système projectif filtrant \deleted{\unsure{de $Y$-pré|de préschémas}} ess. affine de $Y$-préschémas de prés. finie (= de type fini).

(i) Un morphisme affine, un morphisme de prés. finie, sont ess. affine.

(ii)
% archive sheet number: 70
